\section{The end limit and the numerical inequality}
\label{sec:rpi:numerical-limit}

The vanishing deficit determines a limiting radial potential. Its zero
boundary value identifies the largest limiting frontier radius, and
comparison with enclosing spheres gives the sharp full-area bound.

\begin{proposition}[The inequality for a harmonically flat end]
\label{prop:rpi:harmonic-flat-inequality}
Let $(E,h)$ satisfy the hypotheses of Theorem~\ref{thm:rpi}, with
$n\ge8$, nonempty frontier of full area $A$, and a harmonically flat
end. Then
\[
 m_{\rm ADM}(h)\ge\frac12(A/\omega)^{d/(n-1)}.
\]
\end{proposition}

\begin{proof}
The conformal flow preserves $A$ and has decreasing mass
$m(t)\to M\ge0$. Proposition~\ref{prop:rpi:complete-end-factors}
controls its auxiliary factor on a common tail. Put
$q=2(n-1)/d$.

\paragraph{Harmonic normalization and characteristics.}
Recall the total harmonic factors $F_t=h_tU_t$, $Z_t=h_tV_t$,
and $J_t=h_tW_t$ from Section~\ref{sec:rpi:flow-mass-capacity}.
The flow metric has factor $F_t$, and the auxiliary metric has factor
$J_t$; their monopole coefficients are $m(t)/2$ and $\mu(t)/2$.
The background factor $h_t=H\circ\Phi_t$ contributes the initial-mass
term: the $r^{-d}$ coefficient of $U_t$ is $e^{-t}B(t)$, while that
of $F_t$ is $m(t)/2=e^{-t}B(t)+e^{-2t}m(0)/2$.

Set
\[
 f_t=F_t-1-\frac{m(t)}2r^{-d},\qquad
 j_t=J_t-1-\frac{\mu(t)}2r^{-d}.
\]
The evolution and mass identities give
\[
 \partial_t f_t=2(f_t-j_t)+\frac2d r\partial_r f_t.
\]
For $y_s=e^{2(t-s)/d}x$, the chain rule therefore gives
\[
 \frac{d}{ds}\bigl(e^{-2s}f_s(y_s)\bigr)
 =-2e^{-2s}j_s(y_s).
\]
Here is the required quantitative bound for $j_t$.
Equation~\eqref{eq:rpi:auxiliary-weighted-limit} gives
$\sup_{r\ge s_0}r^d|J_t-1|\to0$ for every fixed
$s_0>R_c$. Choose $s_0=2\max\{R_c,R_{\max},1\}$.
The Kelvin transform
\[
 K_t(y)=|y|^{-d}(J_t(y/|y|^2)-1)
\]
extends harmonically across zero, has value
$K_t(0)=\mu(t)/2$, and is uniformly small on
$|y|\le s_0^{-1}$. The interior gradient estimate on the half ball
therefore gives $|j_t(x)|\le C\varepsilon r^{-d-1}$ for all sufficiently
large $t$ and $r\ge2s_0$. For any fixed $0<\beta\le1$, this proves,
for $s\ge t_0$ and $r\ge2s_0$,
\[
 |j_s(x)|\leq C\varepsilon r^{-d-\beta},\qquad
 |f_{t_0}(x)|\leq C_0r^{-d-\beta}.
\]
Integrating the characteristic equation gives the explicit estimate
\[
 |f_t(x)|\leq
 \left[C_0e^{-2\beta(t-t_0)/d}
 +\frac{dC}{\beta}\varepsilon
    \bigl(1-e^{-2\beta(t-t_0)/d}\bigr)\right]r^{-d-\beta}.
\]
The first term tends to zero with the initial time $t_0$ fixed. First taking $t\to\infty$ and
then $\varepsilon\downarrow0$ proves the desired exterior convergence
of the total factors. More explicitly, with $M=\lim m(t)$,
\[
 F_t\longrightarrow1+\tfrac M2r^{-d},\qquad
 Z_t=F_t-2J_t\longrightarrow-1+\tfrac M2r^{-d}.
\]
Since $h_t\to1$ uniformly on any fixed exterior region,
$V_t=Z_t/h_t$ has the latter limit as well.

\paragraph{Boundary control on fixed annuli.}
Choose any sequence $t_i\to\infty$. We obtain a uniform boundary
H\"older estimate away from the collapsing compact core; it will
transfer the radial potential to a farthest limiting frontier point.
Use the direct end coordinates $x=e^{-2t_i/d}y$ on the end chart.
In this coordinate region, denote the rescaled time-$t_i$ exterior
by $E_i$ and its rescaled frontier by $\Sigma_i$.
If the original end chart begins at radius $R_{\mathrm{chart}}$, represent the core
by the closed ball of radius $R_{\mathrm{chart}}e^{-2t_i/d}$. Adjoin that ball to the
rescaled boundary portion in the end and call the resulting compact set
$\mathcal C_i$. The radius bound puts these compact sets in a fixed ball.
They record the frontier and the collapsing core in Euclidean
coordinates; enclosure comparisons will still use the original filled
regions. On every fixed compact annulus
away from the origin, the rescaled upper and lower conformal bounds give
uniform comparison of $G_{t_i}$ with $\delta$. More explicitly, for
an annulus with inner radius $a_0>0$, fix
\[
 L\geq\max\left\{1,\frac d2\log\frac{R_{\max}}{a_0}\right\},
 \qquad S=t_i-L\geq0.
\]
If $|x|\geq a_0$, then $|y|=e^{2t_i/d}|x|\geq R(S)$.
Since $t_i\geq S+1$ and $u_t$ is nonincreasing in time, the radius proof directly gives
\[
 e^{-t_i}\leq u_{t_i}(y)\leq u_{S+1}(y)\leq4e^{-S},
 \qquad 1\leq U_{t_i}(x)\leq4e^L.
\]
Together with $H(e^{2t_i/d}x)\to1$, these are the required uniform
bounds on the total factor $F_{t_i}$ and metric $G_{t_i}$.
The rescaled initial
obstacle is eventually disjoint from that annulus. Thus the enclosure
property permits both filling and deleting local balls and gives a
uniform Euclidean perimeter-quasiminimality ratio. The isoperimetric
argument supplies, with constants independent of $i$, both-phase density
\[
 |B_s(p_i)\cap E_i|\geq\eta s^n,
 \qquad |B_s(p_i)\setminus E_i|\geq\eta s^n,
 \qquad p_i\in\Sigma_i,\quad 0<s<\rho_0.
\]
Here $p_i$ lies in a smaller concentric annulus and $\rho_0$ is chosen so
$B_{4\rho_0}(p_i)$ stays in the annulus with the uniform estimates.
For example, applying the deleting competitor to a phase of volume
$b(s)$ gives $b(s)^{(n-1)/n}\leq Cb'(s)$; integrating this inequality
gives the claimed lower volume bound. Apply the same argument to the
other phase.

The backgrounds $\bar G_{t_i}$ converge smoothly to $\delta$ on the
annulus. Hence the operators for $V_i$ are uniformly elliptic there.
Construct the background equilibrium potential in the original ambient
manifold by minimizing Dirichlet energy subject to value one on the
filled conductor and zero at infinity. Write $\mathcal P_i$ for its
pullback to the end chart. The homogeneous energy construction, or smooth
conductor exhaustion and weak compactness, extends the potential by one
across the conductor. Truncation gives $0\le\mathcal P_i\le1$, and the variational
inequality makes $\mathcal P_i$ a weak supersolution. Thus
$q_i=1-\mathcal P_i$ has a $W^{1,2}_{\mathrm{loc}}$ zero extension, is a weak
subsolution, and satisfies $0\le q_i\le1$. If
$M_i=\sup_{B_{4s}(p_i)}q_i$, the function $M_i-q_i$ is a nonnegative
supersolution and equals $M_i$ on a set of volume at least $\eta s^n$.
The weak Harnack inequality gives
\[
 \sup_{B_s(p_i)}q_i\leq(1-\theta)
       \sup_{B_{4s}(p_i)}q_i,
 \qquad 0<\theta<1.
\]
Iteration yields uniform boundary H\"older control
\[
 |q_i(x)|\leq C\left(\frac{|x-p_i|}{\rho_0}\right)^\kappa
 \quad (|x-p_i|<\rho_0),\qquad \kappa>0.
\]
The variational construction of the flow potential gives
$q_i=-V_i$ on the exterior portion of the end chart; see
Lemma~\ref{lem:rpi:flow-construction} and
\cite[Lemma~3.2]{BiZhu2026}. Thus the estimate controls $|V_i|$
uniformly in $i$ on each fixed annulus, including at singular
frontier points.

\paragraph{The largest limiting radius and the area comparison.}
Take a Hausdorff limit $\mathcal C$ of $\mathcal C_i$, and let $O$ be
the unbounded component of $\mathbb R^n\setminus\mathcal C$.
Every compact subset of $O$ eventually lies in the actual exterior:
join its points to infinity by finitely many paths with positive
distance from $\mathcal C$, then use Hausdorff convergence and the absence of
a boundary crossing along those paths.
The uniform annular bounds and interior elliptic estimates give a
common subsequence of $F_{t_i}$ and $V_i$ converging on compact subsets
of $O$. Their far-exterior limits identify both functions throughout
$O$ by unique continuation:
\[
 F_\infty(x)=1+\frac M2|x|^{-d},\qquad
 V_\infty(x)=-1+\frac M2|x|^{-d},
\]
where $M=\lim_{t\to\infty}m(t)$. The origin lies in $\mathcal C$ and hence
outside the continuation domain $O$.
The boundary H\"older estimate implies
$V_\infty=0$ at every point of $\partial O\setminus\{0\}$.
Set $R_* = \max\{|x|:x\in\mathcal C\}$. If $R_*>0$, choose $p\in\mathcal C$
with $|p|=R_*$. Points on the outward radial ray from $p$ lie in $O$,
so the boundary estimate and the displayed potential give
\[
 0=-1+\frac{M}{2R_*^d},\qquad
 R_*=r_\infty:=(M/2)^{1/d}.
\]
In particular $M>0$, and $\mathcal C\subset\overline B_{r_\infty}$.

If $R_*=0$, then $\mathcal C=\{0\}$ and the same harmonic limit holds on every
fixed annulus. Its sign $V_\infty\leq0$ forces $M=0$.
For any fixed $R>0$, the normalized boundaries are eventually inside
$B_R$, while the total factors converge to $1$ on $S_R$ by harmonic
compactness and unique continuation. The minimizing-enclosure property
would give $A\leq\omega R^{n-1}$ for every $R>0$, contradicting
$A>0$. Thus $R_*>0$ and $M>0$. The enclosing coordinate spheres in
this argument also enclose the original compact core.

Finally, for any $R>r_\infty$, outer containment and the minimizing
property give, after taking the subsequential limit,
\[
 A\leq\omega R^{n-1}
       \left(1+\frac{M}{2R^d}\right)^q.
\]
Letting $R\downarrow r_\infty$ yields
\[
 A\leq\omega(2M)^{(n-1)/d},\qquad
 m(0)\geq M\geq\frac12
        \left(\frac{A}{\omega}\right)^{d/(n-1)}.
\]
\end{proof}
